\chapter{Beyond Gaussianity: Laplace Innovations and the Gaussian-Message
Approximation}\label{ch:laplace}\label{sec:bp-nongaussian}

Every exact result of Chapter~\ref{ch:gaussian} rests on the closure of
the Gaussian family under the two operations the model performs,
multiplying densities and marginalising. This chapter removes that
property in the most controlled way available and answers RQ3. We keep
the AR(1) recursion, the OU channel, and the factor-graph structure, and
change only the innovation law from Gaussian to Laplace: heavy-tailed and
kinked at the origin, yet log-concave and analytically tractable, it is a
minimal step outside the Gaussian family, so every deviation we find is
attributable to non-Gaussianity alone. The chapter first establishes what
remains exact (the score--posterior identity, the tree structure, one
boundary integral, and a differential characterisation of the $K=1$
marginal), then locates precisely where closed forms stop, and finally
measures, against a validated numerical reference, what a Gaussian
message family costs.

\section{What changes, and what survives}\label{sec:l-model}

The Laplace AR(1) chain replaces the Gaussian innovations of
\eqref{eq:ar1} with Laplace ones, keeping the initial law and the
second-order structure fixed:
\begin{equation}
  a_0 \sim \Nd(0,1), \qquad
  a_k = \alpha\, a_{k-1} + \varepsilon_k, \qquad
  \varepsilon_k \sim \mathrm{Laplace}(0, b) \ \text{i.i.d.}, \qquad
  k = 1,\dots,K-1 ,
  \label{eq:l-chain-prior}
\end{equation}
observed through the same per-frame OU channel. The innovation scale is
\emph{not} free: $b = \sqrt{(1-\alpha^2)/2}$, so that
$\operatorname{Var}(\varepsilon_k) = 1-\alpha^2$ and the chain has unit
marginal variance and $\operatorname{Cov}(a_j,a_k)=\alpha^{|j-k|}$, exactly
as in the Gaussian case. This is what makes the comparison in this chapter
a comparison: the Gaussian and Laplace chains agree on every second moment
and differ \emph{only} beyond them, the innovation's excess kurtosis
($0$ against $+3$) being the single non-Gaussianity knob. The same
convention --- Gaussian initial state, variance-matched innovations --- is
used for every non-Gaussian family in this thesis, and carries forward
unchanged to the estimated chain of \eqref{eq:em-chain}. Three structural facts
survive unchanged, because none of them used Gaussianity of the prior.
The score--posterior identity \eqref{eq:tweedie-joint} holds for any
$p_0$. The posterior factorises over the same tree-shaped factor graph
(Proposition~\ref{prop:chain-tree}), so the sum--product recursion
\eqref{eq:bp-fwd}--\eqref{eq:bp-bwd} still computes exact smoothing
marginals in one forward and one backward sweep. And the OU likelihood
factors are
still Gaussian in $a_k$. What is lost is representability: the messages
are no longer Gaussian, no finite-dimensional family we know of is closed
under the updates, and the cost of an update is no longer bounded. The
score also ceases to be linear in $x$, which is the first thing to make
precise.

\section{One frame: the score becomes nonlinear}\label{sec:l-K1}

For $K = 1$ there is no chain; the model isolates the ingredient
``non-Gaussian prior meets Gaussian channel''. Let
$a \sim \mathrm{Laplace}(0, b)$, $p_0(a) = e^{-|a|/b}/2b$, and write
$\mu := e^{-t}$. The noised marginal is the convolution
\begin{equation}
  p_t^{(1)}(x) = \int p_0(a)\, \Nd\big(x;\, \mu a,\, \Deltat\big)\, \dd a .
  \label{eq:l-pt1}
\end{equation}
As $t \to 0$ the Gaussian kernel concentrates and
$p_t^{(1)} \to p_0$; as $t \to \infty$, $p_t^{(1)} \to \Nd(0,1)$, the OU
attractor. By the score--posterior identity,
$S(x,t) = (\mu\,\E[a|x] - x)/\Deltat$; for a Laplace prior the
conditional mean is a nonlinear function of $x$, and the score inherits
the nonlinearity. Its limits are exact: for every fixed $x \neq 0$,
\begin{equation}
  t \to 0:\quad S(x, t) \;\longrightarrow\;
  \partial_x \log p_0(x) = -\frac{\operatorname{sign}(x)}{b},
  \qquad\qquad
  t \to \infty:\quad S(x, t) \;\longrightarrow\; -x .
  \label{eq:l-score-limits}
\end{equation}
At small $t$ the score converges pointwise to the sign-shaped score of
the prior, with the nonlinearity concentrated in the kink near $x = 0$;
at large $t$ it flattens into the linear attractor score $-x$
(Figure~\ref{fig:l-k1}). The audit verifies the posterior-expectation
route against central finite differences of $\log p_t^{(1)}$ computed by
quadrature: maximum disagreement $1.19 \times 10^{-7}$ over a grid of
$(b, t, x)$, at the truncation floor of the finite difference.

\begin{figure}[t]\centering
\includegraphics[width=\linewidth]{fig_laplace_k1.pdf}
\caption{The Laplace $K=1$ problem ($b = 1$). Left: noised marginal
$p_t^{(1)}$ at increasing diffusion times, from the kinked Laplace prior
(dotted) to the Gaussian attractor (dashed). Right: the score at the same
times; the dashed line is the attractor $-x$. The nonlinearity lives near
$x = 0$ and flattens as $t$ grows; the limits are the exact statements
\eqref{eq:l-score-limits}. Curves computed by deterministic quadrature.}
\label{fig:l-k1}
\end{figure}

\subsection{A differential characterisation of the noised marginal}
\label{sec:l-fourier}

The channel writes $x$ as a sum of independent variables, so
\eqref{eq:l-pt1} is a convolution, which factorises in Fourier space:
\begin{equation}
  \widehat{p}_t^{\,(1)}(k)
  = \frac{1}{1 + b^2 \mu^2 k^2}\; e^{-\Deltat k^2/2},
  \label{eq:l-cf}
\end{equation}
the product of the (rescaled) Laplace and Gaussian characteristic
functions.

\begin{theorem}[Screened Poisson equation]\label{thm:l-screened}
The Laplace noised marginal solves the linear, constant-coefficient ODE
\begin{equation}
  \boxed{\;
  \big( 1 - b^2 \mu^2\, \partial_x^2 \big)\, p_t^{(1)}(x)
  \;=\; G_{\Deltat}(x),
  \;}
  \qquad
  G_{\Deltat}(x) = \frac{e^{-x^2/2\Deltat}}{\sqrt{2\pi \Deltat}} .
  \label{eq:l-screened}
\end{equation}
\end{theorem}

\begin{toolbox}{From the characteristic function to the screened Poisson
equation}
Multiply \eqref{eq:l-cf} by $(1 + b^2\mu^2 k^2)$:
\begin{equation*}
  \big( 1 + b^2 \mu^2 k^2 \big)\, \widehat{p}_t^{\,(1)}(k)
  = e^{-\Deltat k^2 / 2},
\end{equation*}
whose right-hand side is the characteristic function of $G_{\Deltat}$.
Multiplication by $k^2$ in Fourier space is $-\partial_x^2$ in real
space, so the polynomial $(1 + b^2\mu^2 k^2)$ becomes the operator
$(1 - b^2\mu^2 \partial_x^2)$, and \eqref{eq:l-screened} follows by
inverse transformation. The audit checks \eqref{eq:l-cf} against direct
quadrature of the characteristic function to $2.9 \times 10^{-6}$, the
floor of oscillatory quadrature. \hfill$\square$
\end{toolbox}

Equation \eqref{eq:l-screened} states that the noised marginal is the
solution of a screened Poisson equation with a Gaussian source, with
screening length $b\mu = b e^{-t}$: the memory of the prior's tails
shrinks as $e^{-t}$. The deviation from Gaussianity is governed by a
single local differential operator whose strength dies with diffusion
time.

One further exact consequence identifies what a Gaussian approximation
must get wrong. Differentiating the score--posterior identity once more,
\begin{equation}
  H_t(x) := -\partial_x^2 \log p_t^{(1)}(x)
  = \frac{1}{\Deltat}
    \left( 1 - \frac{\mu^2}{\Deltat}\, \operatorname{Var}[a \,|\, x] \right).
  \label{eq:l-curvature}
\end{equation}
For a Gaussian prior the posterior variance is independent of $x$ and
$H_t$ is constant. For the Laplace prior $\operatorname{Var}[a|x]$ depends
on $x$ and $H_t$ is a nontrivial function: the breakdown of a constant
curvature is the second-order form of the breakdown of score linearity.
A Gaussian message is exactly the object that cannot represent an
$x$-dependent curvature, so \eqref{eq:l-curvature} states what the
Gaussian closure of Section~\ref{sec:l-closure} will miss, before it is
measured.

\section{The chain: where closed forms stop}\label{sec:l-chain}

\begin{theorem}[Joint score from chain messages]\label{thm:l-markov-score}
For any chain prior, in particular \eqref{eq:l-chain-prior}, and the OU
likelihood, the joint score is, componentwise,
\begin{equation}
  S_k(x, t) = \frac{\mu\, \E[a_k | x] - x_k}{\Deltat},
  \qquad
  \E[a_k | x]
  = \frac{\displaystyle\int a_k\; \mu_{\to k}(a_k)\;
      \psi^{(k)}_{\mathrm{ob}}(a_k)\;
      \mu_{\leftarrow k}(a_k)\;\dd a_k}
     {\displaystyle\int \mu_{\to k}(a_k)\;
      \psi^{(k)}_{\mathrm{ob}}(a_k)\;
      \mu_{\leftarrow k}(a_k)\;\dd a_k},
  \label{eq:l-markov-score}
\end{equation}
with the messages of Proposition~\ref{prop:bp-convA}.
\end{theorem}

This is a complexity statement: the $K$-dimensional posterior integral is
replaced by two one-dimensional message recursions plus one
one-dimensional integral per frame. Computational locality survives
non-Gaussianity; what does not survive is the collapse of those
one-dimensional integrals into closed-form arithmetic. Exactly one such
integral has a closed form.

\begin{proposition}[One-step boundary integral]\label{prop:l-onestep}
For the Laplace transition $M(a'|a) = \frac{1}{2b} e^{-|a' - \alpha a|/b}$
and OU evidence $\Nd(x'; \mu a', \Deltat)$,
\begin{equation}
  \int_\R M(a' | a)\; \Nd(x'; \mu a', \Deltat)\; \dd a'
  \;=\; p_t^{(1)}\big( x' - \mu\alpha\, a \big) :
  \label{eq:l-onestep}
\end{equation}
the first backward message of the chain is the $K = 1$ noised marginal
evaluated at a shifted argument.
\end{proposition}

The proof is the innovation substitution $r := a' - \alpha a$, which
turns the transition into the bare Laplace density at $r$ and the
evidence into $\Nd(x' - \mu\alpha a;\, \mu r,\, \Deltat)$; the integral
is then the definition \eqref{eq:l-pt1}. The shift $\mu\alpha a$ is what
frame $a$ predicts for the next noisy observation, and the whole $K = 1$
analysis, screened Poisson equation included, recycles as the chain's
boundary building block.

\begin{remark}[Why iteration destroys the closed form]\label{rmk:l-breakdown}
Equation \eqref{eq:l-onestep} works because the integrand contains
exactly one Laplace kink times exactly one Gaussian. Every subsequent BP
step violates this: the incoming message already carries the propagated
structure of several frames, so the next integrand contains two or more
absolute-value kinks. On each sign-quadrant the integral reduces to
half-line Gaussians, but globally it is a sum of error functions whose
arguments depend on the conditioning variable non-affinely, and the clean
form does not survive one iteration. In the Gaussian chain every step is
Gaussian-times-Gaussian and closure propagates forever; this is the
precise structural location of the difference between the two models, and
the reason the Gaussian projection below is an approximation with an
error to be measured, not an identity.
\end{remark}

A change of coordinates shows the obstruction is intrinsic, not an
artefact of the frame basis. In innovation coordinates $r_0 = a_0$,
$r_k = a_k - \alpha a_{k-1}$ (unit Jacobian), the prior factorises into
independent variables --- $r_0$ Gaussian and $r_{k\geq1}$ i.i.d.\
Laplace --- but the likelihood becomes
$p(x|r) \propto \exp(-\tfrac{\mu^2}{2\Deltat} r^\T L^\T L r
+ \tfrac{\mu}{\Deltat} x^\T L r)$ with $L$ the unit lower-triangular
map $a = Lr$, and $L^\T L$ is dense. The coupling does not disappear; it
migrates from the prior to the likelihood. There is no basis in which
both are diagonal at once.

\section{A validated numerical reference: grid BP}\label{sec:l-grid}

Since functional messages admit no closed form, the exact recursion is
approximated numerically: messages are represented by their values on $M$
equispaced points in $[-A, A]$ and the update integrals are carried out
by trapezoidal quadrature, each update touching every pair of grid
points. This
\emph{grid BP} is a numerical approximation of the functional-message
recursion, and it can serve as a reference only inside a validated
regime, which we establish on the exactly solvable Gaussian chain. There
the trapezoid rule on analytic, rapidly decaying integrands converges
spectrally, so failure is abrupt rather than gradual: at $M = 101$,
$A = 8$ the score error sits at the $10^{-15}$ floor down to $t = 0.02$
and fails only when the likelihood width $\sqrt{2t}$ approaches the grid
step. The working configuration for all non-Gaussian experiments is
$M = 401$, $A = 8$, with a reference solution at $M = 801$ and the
resolution constraint $\sqrt{2t} \gtrsim 3\,\dd x$ checked automatically;
grid self-convergence error remains $10^4$ to $10^7$ times smaller than
every effect measured below. Within this regime we treat grid BP as the
reference score; outside it no claim is made.

All score errors are amplified posterior-mean errors: for any two solvers
with means $\widehat m$ and $m_{\mathrm{ref}}$, the score--posterior
identity gives, exactly,
\begin{equation}
  \widehat s - s_{\mathrm{ref}}
  = \frac{e^{-t}}{\Deltat}\,
    \big( \widehat m - m_{\mathrm{ref}} \big),
  \qquad
  \frac{e^{-t}}{\Deltat} \sim \frac{1}{2t} \;\; (t \downarrow 0).
  \label{eq:bp-master-error}
\end{equation}
The numerical residual of this identity is logged in every experiment and
stays at $10^{-15}$ to $10^{-12}$. Equation \eqref{eq:bp-master-error}
explains in advance why small diffusion times are the difficult regime:
whatever an approximation does to the posterior mean is amplified by
$1/2t$ in the score.

\section{The Gaussian-message approximation, measured}\label{sec:l-closure}

The object under test is Gaussian-projected BP: the same sweeps, with
every message forced into the Gaussian family by moment matching. This is
the algorithm one would deploy when quadrature is unaffordable, and the
chain analogue of the per-node closures of statistical physics
(Appendix~\ref{app:amp}).

On the Laplace chain (chain correlation $\alpha = 0.85$, $K = 40$ frames, $60$ trials per
$t$, reference $M = 801$), the median relative score error of the
Gaussian closure is
\begin{center}
\begin{tabular}{@{}lccccc@{}}
\toprule
$t$ & $0.05$ & $0.32$ & $0.98$ & $2.07$ & $3.0$ \\
\midrule
median relative score error & $0.24$ & $4.6\times10^{-2}$
& $6.8\times10^{-3}$ & $2.5\times10^{-4}$ & $7.3\times10^{-6}$ \\
$\Prob(\mathrm{err} > 0.1)$ & $1.00$ & $0.08$ & $0.00$ & $0.00$ & $0.00$ \\
\bottomrule
\end{tabular}
\end{center}
\noindent(Figure~\ref{fig:bp-closure}). Two thresholds are worth separating,
because they are different statistics and sit an order of magnitude apart in
$t$. The probability of a trial exceeding $10\%$ relative error first reaches
zero at $t \approx 0.98$; the \emph{median} error first falls below $10^{-3}$
only at $t \approx 2.1$. In between --- around $t \approx 1$, where the
failure probability has already vanished --- the typical error is still
$\sim\!7\times10^{-3}$.

The closure error is of order tens of percent at small diffusion time and
decays as diffusion noise Gaussianises the tilted posteriors. Inspection of
the worst trials shows the mechanism: the reference belief is sharply
non-Gaussian, occasionally bimodal; the single Gaussian commits to a
compromise mean; and \eqref{eq:bp-master-error} amplifies the shift. The
score \emph{direction} is much more robust than its magnitude: the median
cosine similarity between projected and reference scores is $0.976$ at
$t = 0.05$ (where the magnitude error is $24\%$), exceeds $0.99$ for all
$t \geq 0.11$, and the worst single trial of the sweep stays above $0.879$.
These are numerical observations for the stated parameter ranges, not
theorems.

\begin{figure}[t]\centering
\includegraphics[width=0.74\linewidth]{fig_laplace_closure.pdf}
\caption{Gaussian-closure score error on the Laplace chain
($\alpha = 0.85$, $K = 32$, $60$ trials per $t$; median and 10--90\%
band) against the grid error budget measured on identical data. The
closure error exceeds the reference error by $10^4$ to $10^7$ at every
$t$, so the measurement is trustworthy where it is used.}
\label{fig:bp-closure}
\end{figure}

\subsection{Varying the innovation law}

Sweeping six innovation families with identical covariance
$\alpha^{|i-j|}$ (Laplace, Student-$t$ with $\nu = 5, 8$, and symmetric
Gaussian mixtures of increasing separation) isolates the shape of the
innovation law as the only moving part. Three observations, all within
the tested grids. First, the noise level dominates: every family's error
decays to about $10^{-3}$ or below by $t = 2.4$; all substantial closure
error lives at $t \lesssim 0.4$. Second, bimodality is the worst case: at
$t = 0.05$, $\alpha = 0.5$, the median error is $0.94$ for a strongly
bimodal mixture, $0.52$ for Laplace, $0.32$ for Student-$t_{\nu=5}$, and
$0.07$ for a mildly bimodal mixture; a single Gaussian cannot represent
two posterior modes and commits to the gap between them. Third, stronger
chain correlation helps rather than hurts: for Laplace at $t = 0.05$ the
error is $0.43$, $0.27$, $0.13$ at $\alpha = 0.5$, $0.85$, $0.95$;
informative neighbours act like additional Gaussian observations and
Gaussianise the local tilted posteriors.

\section{What can and cannot be concluded}\label{sec:l-open}

Exact: the surviving structure of Section~\ref{sec:l-model}, the limits
\eqref{eq:l-score-limits}, the factorisation \eqref{eq:l-cf}, the
screened Poisson equation \eqref{eq:l-screened}, the curvature identity
\eqref{eq:l-curvature}, the message form \eqref{eq:l-markov-score}, the
boundary integral \eqref{eq:l-onestep}, the breakdown argument of
Remark~\ref{rmk:l-breakdown}, and the error identity
\eqref{eq:bp-master-error}. Numerical, within validated regimes: the
quadrature audits of Section~\ref{sec:l-K1}, the grid-BP calibration, and
the closure-error measurements, which are specific to the stated
innovation laws, correlation levels, chain lengths, and diffusion times.
Not concluded: closed forms for general (non-boundary) messages at
$K \geq 2$; any locality law for the non-Gaussian score beyond the
structural delimitation above; and any statement about learned scores. An
earlier conjecture of this project, that the Laplace $K = 2$ posterior
Hessian is approximately rank-two, proved regime-dependent under
numerical tests and is not asserted.
